\chapter{Structure du marché des options}\label{ch:m1:options-market-structure}

Un client de détail demande à «~acheter 10 calls au marché~». L'ordre n'atteint
jamais un carnet d'ordres au sens ordinaire. Son courtier l'envoie à un
\omterm{def:m1:retail-flow-and-wholesaling:wholesaler}{wholesaler}~; le \omterm{def:m1:retail-flow-and-wholesaling:wholesaler}{wholesaler} le soumet à une bourse avec son propre ordre pour prendre
la contrepartie à un prix garanti~; la bourse diffuse une demande de réponses et
attend, pendant une durée qui peut n'être que d'un dixième de seconde, tandis que
les plus grandes sociétés de \omterm{def:m1:what-a-trading-firm-does:market-maker}{teneur de marché} du monde décident de proposer un centime de mieux.
Puis l'ordre est réparti~: une partie revient à ceux qui ont amélioré le prix, une part réservée
à la société qui l'a apporté, le reste étant partagé au prorata de la taille.
Chacune de ces étapes est inscrite dans une règle de la bourse. Le marché des options
est le lieu où les mécanismes des cinq derniers chapitres (wholesaling, \omterm{def:m1:matching-algorithms-and-implied-spreads:prorata}{allocation au prorata},
privilèges des \omterm{def:m1:what-a-trading-firm-does:market-maker}{teneurs de marché}, modèles de frais) se combinent le plus densément,
et où la position d'une société dans la structure décide davantage de son profit
que ses modèles.

\section{Un marché conduit par les cotations}

\begin{definition}[Teneur de marché d'options et obligation de cotation]\label{def:m1:options-market-structure:mm}
Un \emph{teneur de marché d'options}\index{teneur de marché d'options} est une
société enregistrée auprès d'une bourse pour coter des classes d'options. Elle
accepte une \emph{obligation de cotation}\index{obligation de cotation}, en
général de coter les deux côtés dans un pourcentage déterminé des séries de ses
classes, pendant un pourcentage déterminé de la séance de négociation, à
l'intérieur d'une largeur maximale, et reçoit en contrepartie des privilèges~:
frais réduits, droits d'allocation, droit d'envoyer des cotations en masse.
\end{definition}

Avec deux mille séries sur une même action (\cref{ch:m1:option-contracts-and-exchanges}),
dont la plupart ne se négocient que quelques fois par jour, voire pas du tout, les acheteurs
et les vendeurs naturels ne se rencontrent presque jamais. Presque chaque prix affiché
est une cotation de \omterm{def:m1:what-a-trading-firm-does:market-maker}{teneur de marché}, calculée à partir du cours de l'action et d'une
surface de volatilité, et actualisée dès que l'un des deux bouge. Le carnet est une
projection d'un modèle, et les questions économiques sont de savoir qui peut coter,
ce qu'il doit coter, et comment un ordre entrant est réparti entre eux.

\begin{definition}[Protection des cotations]\label{def:m1:options-market-structure:protection}
La \emph{protection des cotations}\index{protection des cotations} désigne les
mécanismes fournis par la bourse qui limitent l'exposition d'un \omterm{def:m1:what-a-trading-firm-does:market-maker}{teneur de marché} sur
l'ensemble de ses cotations~: un moniteur de risque qui annule toutes les cotations d'une société
dans une classe lorsque ses exécutions, sur une courte fenêtre, dépassent un nombre fixé,
des messages d'annulation en masse et des coupe-circuits.
\end{definition}

Un \omterm{def:m1:what-a-trading-firm-does:market-maker}{teneur de marché} qui cote deux mille séries offre, économiquement, un seul pari
sur l'action sous deux mille formes. Un trader informé qui prend
cent d'entre elles en une milliseconde a échangé cent fois contre la même opinion périmée.
Sans protection, les cotations ne vaudraient qu'une fraction de leur taille.

\section{Allocation~: les clients d'abord, puis la taille}

\begin{definition}[Priorité client]\label{def:m1:options-market-structure:customer}
Sous la \emph{priorité client}\index{priorité client}, les ordres au repos des
\emph{clients prioritaires} (au sens large, le public, par opposition aux \omterm{def:m1:the-sell-side:sell-side}{broker-dealers}
et aux traders professionnels) au meilleur prix sont exécutés avant tout
intérêt professionnel à ce prix, quelle que soit leur heure d'arrivée. Le
reste est réparti entre les professionnels, en général au prorata de la taille,
souvent après un droit accordé à un \omterm{def:m1:what-a-trading-firm-does:market-maker}{teneur de marché} désigné.
\end{definition}

\begin{omfigure}
\begin{tikzpicture}
\begin{axis}[width=11.5cm,height=5.2cm,ybar,bar width=8pt,
  ylabel={lots reçus},
  symbolic x coords={MM1,C1,DMM,F,C2,MM2},xtick=data,
  xticklabels={{MM1~: 100},{C1~: 10},{DMM~: 100},{F~: 50},{C2~: 5},{MM2~: 250}},
  tick label style={font=\small},label style={font=\small},xticklabel style={font=\footnotesize\strut},
  ymin=0,ymax=66,ymajorgrids,grid style={gray!25},enlarge x limits=0.1,
  nodes near coords,nodes near coords style={font=\scriptsize},
  legend columns=3,legend style={font=\footnotesize,at={(0.5,-0.2)},anchor=north,draw=none,fill=none,/tikz/every even column/.append style={column sep=8pt}},
  table/col sep=comma]
\addplot[fill=omExo!50,draw=omExo] table[x=order,y=pro_rata]{figdata/markets-1/24-options-market-structure/rules.csv};
\addplot[fill=omDef!70,draw=omDef] table[x=order,y=customer]{figdata/markets-1/24-options-market-structure/rules.csv};
\addplot[fill=omThm!60,draw=omThm] table[x=order,y=entitled]{figdata/markets-1/24-options-market-structure/rules.csv};
\legend{allocation au prorata simple,priorité client,priorité client et un droit de 40~\%}
\end{axis}
\end{tikzpicture}

\omcaption{Un ordre de 115 lots rencontre un niveau de six ordres au repos, classés
par priorité temporelle avec leurs tailles~: deux clients (C), trois \omterm{def:m1:what-a-trading-firm-does:market-maker}{teneurs de
marché}, dont un désigné (DMM), et une société (F). Les clients sont exécutés
intégralement sous la priorité~; le droit double la part du \omterm{def:m1:what-a-trading-firm-does:market-maker}{teneur de marché}
désigné au détriment des autres professionnels. Données~: la
construction du chapitre.}\label{fig:m1:options-market-structure:rules}
\end{omfigure}

La règle est un transfert. Les clients, dont les ordres portent le moins
d'information, valent le plus à négocier contre eux~; leur accorder la priorité les attire
vers la bourse, et leur flux devient alors l'enjeu que se disputent les règles
d'allocation des professionnels. Les bourses qui veulent attirer les cotations des
\omterm{def:m1:what-a-trading-firm-does:market-maker}{teneurs de marché} offrent des droits et le prorata~; celles qui veulent attirer
les preneurs de liquidité offrent la \omterm{def:m1:matching-algorithms-and-implied-spreads:fifo}{priorité prix-temps} et des rabais. La même série
se négocie sous les deux régimes à la fois, sur des bourses différentes d'un même groupe.

\section{Enchères d'amélioration de prix}

\begin{definition}[Enchère d'amélioration de prix]\label{def:m1:options-market-structure:auction}
Dans une \emph{enchère d'amélioration de prix}\index{enchère d'amélioration de prix}, un
membre soumet l'ordre d'un client (l'ordre d'agence) avec son propre ordre
opposé garantissant l'exécution de l'ensemble à un prix stop égal ou meilleur
que le meilleur cours affiché. La bourse annonce l'enchère, recueille les
réponses pendant une courte période fixe, et exécute l'ordre d'agence contre les
réponses les mieux cotées en premier. Au prix final, le membre initiateur conserve
une part garantie.
\end{definition}

\begin{dated}{2026-09}{L'enchère d'une bourse, d'après ses dépôts de règles}\label{dat:m1:options-market-structure:aim}
La période d'enchère est fixée par classe, d'au moins 100 millisecondes et
d'au plus 3 secondes. Si l'enchère ne produit pas de meilleur prix, les clients
prioritaires au repos dans le carnet au prix stop sont exécutés en premier~; l'ordre
initiateur reçoit ensuite jusqu'à 50~\% de l'ordre d'agence si un autre
participant se trouve à ce prix, ou 40~\% s'il y en a deux ou plus~; le reste est
partagé entre les répondants. La part de l'initiateur au prix final ne peut jamais
dépasser ces pourcentages de l'ordre initial.
\end{dated}

\begin{omfigure}
\begin{tikzpicture}[font=\small]
\draw[thick,-{Stealth[length=2.5mm]}] (0,0) -- (13,0);
\foreach \x/\d/\t in {0.8/$t_0$/{les ordres appariés\\arrivent},4.4/{$t_0$ plus quelques ms}/{l'enchère\\est diffusée},11.2/{$t_0 + 100$ ms}/allocation}{
  \draw[thick] (\x,0.12) -- (\x,-0.12);
  \node[above=3pt,font=\footnotesize\bfseries] at (\x,0.12) {\d};
  \node[below=3pt,font=\footnotesize,align=center] at (\x,-0.12) {\t};}
\foreach \x in {6.2,7.0,7.3,8.6,9.5}{\draw[omThm,thick,-{Stealth[length=1.6mm]}] (\x,0.95) -- (\x,0.08);}
\node[font=\footnotesize,text=omThm] at (7.85,1.2) {réponses, chacune à son propre prix};
\end{tikzpicture}

\omcaption{Une \omterm{def:m1:options-market-structure:auction}{enchère d'amélioration de prix} à la période minimale. Un répondant
dispose d'un dixième de seconde, diminué de sa distance à la bourse, pour recevoir le
message, valoriser l'option et répondre.}\label{fig:m1:options-market-structure:timeline}
\end{omfigure}

L'enchère donne au client un prix non moins favorable que le stop, parfois
meilleur~; elle donne à la firme initiatrice une part garantie du flux qu'elle a choisi~;
et elle laisse à tous les autres la possibilité de prendre l'ordre en améliorant le prix. Son
caractère compétitif est une question empirique, qui dépend de deux grandeurs~: le nombre de
firmes qui répondent et la fréquence à laquelle elles améliorent le prix
(\cref{fig:m1:options-market-structure:auction}).

\begin{omfigure}
\begin{tikzpicture}
\begin{axis}[width=11cm,height=5cm,
  xlabel={probabilité qu'un répondant améliore d'un tick},ylabel={part de l'initiateur (\%)},
  tick label style={font=\small},label style={font=\small},
  xmin=0,xmax=1,ymin=0,ymax=75,grid=major,grid style={gray!25},
  legend columns=2,legend style={font=\footnotesize,at={(0.5,-0.3)},anchor=north,draw=none,fill=none,/tikz/every even column/.append style={column sep=8pt}},
  table/col sep=comma]
\addplot[omDef,very thick,mark=*,mark size=1.4pt] table[x=p_improve,y=share_1]{figdata/markets-1/24-options-market-structure/auction.csv};
\addplot[omThm,very thick,dashed,mark=square*,mark size=1.4pt] table[x=p_improve,y=share_3]{figdata/markets-1/24-options-market-structure/auction.csv};
\legend{un répondant en moyenne,trois répondants en moyenne}
\end{axis}
\end{tikzpicture}

\omcaption{Enchères simulées portant sur un ordre de 100 lots~: part moyenne de la
firme initiatrice. Avec un répondant en moyenne, qui n'améliore jamais, elle conserve
deux tiers~; avec trois répondants qui améliorent la moitié du temps, un huitième. La
garantie ne vaut que ce que la concurrence lui laisse valoir. Données~: la
simulation du tutoriel.}\label{fig:m1:options-market-structure:auction}
\end{omfigure}

\section{Ordres complexes}

\begin{definition}[Carnet d'ordres complexes]\label{def:m1:options-market-structure:complex}
Un \emph{carnet d'ordres complexes}\index{carnet d'ordres complexes} est un carnet destiné aux
stratégies à plusieurs jambes (spreads, straddles, combinaisons avec l'action), cotées à un
prix net unique et exécutées comme un bloc. Un ordre complexe peut être exécuté contre
d'autres ordres complexes ou, jambe par jambe, contre les carnets simples, la bourse
veillant à ce que chaque jambe s'exécute à un prix admis.
\end{definition}

C'est l'homologue, côté options, du \omterm{def:m1:matching-algorithms-and-implied-spreads:calendar}{calendar spread} et de ses ordres implicites
(\cref{def:m1:matching-algorithms-and-implied-spreads:implied})~: un spread vertical dont
chaque jambe est large de dix cents peut être coté et négocié à un cent ou deux de sa
fair value en tant que paquet, car le paquet ne porte qu'une fraction du risque de
chacune des jambes. Une part importante du volume institutionnel d'options passe par de
tels paquets, et un \omterm{def:m1:what-a-trading-firm-does:market-maker}{teneur de marché} qui ne cote que des séries individuelles ne voit
qu'une partie du flux.

\section{Incréments, frais et flux de cotations}

\begin{definition}[Penny program]\label{def:m1:options-market-structure:penny}
Le \emph{penny program}\index{penny program} est la règle, commune à l'ensemble du secteur,
qui fixe l'incrément minimal de cotation des options cotées par classe~: un cent pour
les classes les plus actives, cinq et dix cents pour les autres.
\end{definition}

\begin{definition}[Marketing fee]\label{def:m1:options-market-structure:marketing}
Un \emph{marketing fee}\index{marketing fee} est une redevance prélevée par une bourse sur
les transactions des \omterm{def:m1:what-a-trading-firm-does:market-maker}{teneurs de marché} avec des clients, et mise en commun pour rémunérer les
courtiers qui acheminent les ordres clients vers cette bourse~: un \omterm{def:m1:retail-flow-and-wholesaling:pfof}{paiement pour flux
d'ordres} collecté et redistribué par la bourse elle-même.
\end{definition}

\begin{definition}[OPRA]\label{def:m1:options-market-structure:opra}
\emph{OPRA}\index{OPRA}, l'Options Price Reporting Authority, est le
\omterm{def:m1:us-equity-market-structure:sip}{processeur d'information sur les titres} des options cotées aux États-Unis~: il consolide les
cotations et les transactions de toutes les bourses d'options en un flux unique et publie les
\omterm{def:m1:us-equity-market-structure:nbbo}{meilleurs cours acheteur et vendeur nationaux} de chaque série.
\end{definition}

\begin{dated}{2026-09}{Incréments et messages}\label{dat:m1:options-market-structure:penny}
\emph{Incréments.} Le penny pilot est devenu en 2020 le Penny Interval Program
permanent. Toutes les séries de trois des fonds d'options les plus actifs se cotent en cents
à tout prix~; les autres classes du programme se cotent en cents en dessous de 3.00~\$ et
par paliers de cinq cents à partir de 3.00~\$~; une classe reste dans le programme tant
qu'elle figure parmi les 425 plus actives. \emph{Le flux.} L'avis de capacité du processeur,
publié en septembre 2025, prévoit pour juillet 2026 311~milliards de messages par jour et
un pic de 13.6~millions de messages par 100~millisecondes sur un flux, soit 4.4~gigabits
pendant ce dixième de seconde, à doubler pour le flux redondant et à majorer d'un
dixième pour les retransmissions~; la latence médiane à travers le processeur est inférieure
à 18~microsecondes.
\end{dated}

\begin{omfigure}
\begin{tikzpicture}
\begin{axis}[width=9.5cm,height=4.4cm,ybar,bar width=14pt,
  ylabel={milliards par jour},
  symbolic x coords={10/2025,1/2026,7/2026,1/2027,7/2027},xtick=data,
  tick label style={font=\small},label style={font=\small},xticklabel style={font=\small\strut},
  ymin=0,ymax=400,ymajorgrids,grid style={gray!25},
  nodes near coords,nodes near coords style={font=\footnotesize,fill=white,inner sep=1pt},table/col sep=comma]
\addplot[fill=omDef!60,draw=omDef] table[x=date,y=messages_bn_day]{figdata/markets-1/24-options-market-structure/opra.csv};
\end{axis}
\end{tikzpicture}

\omcaption{Capacité projetée du flux consolidé d'options, en milliards de messages par
jour, d'après l'avis du processeur de septembre 2025. Une firme qui
veut toutes les cotations de toutes les séries doit pouvoir la lire. Données~:
l'avis.}\label{fig:m1:options-market-structure:opra}
\end{omfigure}

Un tick de cinq cents sur une option à trois dollars représente 1.7~\% de son prix~:
l'incrément constitue un plancher sous la fourchette, une source de revenus pour les \omterm{def:m1:what-a-trading-firm-does:market-maker}{teneurs de marché},
et la raison pour laquelle les \omterm{def:m1:options-market-structure:auction}{enchères d'amélioration de prix} ont quelque chose à améliorer. Le volet des
frais est aussi varié qu'en actions (\cref{ch:m1:exchanges-brokers-venues}),
avec une dimension supplémentaire~: les frais dépendent de \emph{qui} négocie (client,
client professionnel, firme, \omterm{def:m1:what-a-trading-firm-does:market-maker}{teneur de marché}) autant que de savoir si l'ordre
apporte ou retire de la liquidité. Le prix net d'un \omterm{def:m1:what-a-trading-firm-does:market-maker}{teneur de marché} sur une transaction n'est connu qu'une
fois la capacité de la contrepartie établie.

\section{Tutoriel~: répartition d'un ordre}\label{tut:m1:options-market-structure:alloc}

\begin{tutorial}
\textbf{Objectif.} Allouer un ordre entrant selon la \omterm{def:m1:options-market-structure:customer}{priorité client} avec un
droit de \omterm{def:m1:what-a-trading-firm-does:market-maker}{teneur de marché} désigné~; lancer une \omterm{def:m1:options-market-structure:auction}{enchère d'amélioration de prix}~; mesurer comment
la concurrence érode la garantie de l'initiateur. \textbf{État final~:} les trois
figures de données de ce chapitre.
\begin{tutsteps}
\item \textbf{Au niveau de la cotation.} Les clients, puis le \omterm{def:m1:what-a-trading-firm-does:market-maker}{teneur de marché} désigné, puis les
autres professionnels par taille.
\omcode{code/firm/match/firm_optmatch.py}{26}{51}{Priorité client, droit acquis et prorata pour le reste.}\label{lst:m1:options-market-structure:alloc}
\item \textbf{Dans l'enchère.} Les meilleurs prix d'abord~; au stop, les clients,
puis la part de l'initiateur, puis les répondants~; l'initiateur prend ce qui
reste, puisqu'il a garanti l'ordre.
\omcode{code/firm/match/firm_optmatch.py}{96}{120}{La fin d'une enchère d'amélioration de prix~: ce qui se passe au prix stop, après que les réponses à meilleur prix ont été exécutées.}\label{lst:m1:options-market-structure:auction}
\item \textbf{Nombreuses enchères.} Des répondants suivant un processus de Poisson, chacun améliorant avec une
probabilité donnée.
\end{tutsteps}
\textbf{Pour aller plus loin.} Laisser les répondants enchérir sur une partie de l'ordre, en
proportion de leur confiance, et laisser l'initiateur choisir un prix stop
à l'intérieur de la fourchette~: trouver le stop qui maximise son profit espéré lorsque son
avantage est de six cents au prix affiché.
\end{tutorial}

\section{Construction~: allocation des options}\label{bld:m1:options-market-structure:optmatch}

\begin{build}
\bfield{Rôle} La plateforme d'options de la firme miniature, dans sa bourse simulée, a besoin des deux
mécanismes qui la distinguent de la plateforme de contrats à terme de
la \cref{bld:m1:matching-algorithms-and-implied-spreads:match}.
\bfield{Interface} Une \texttt{Quote(oid, qty, capacity, designated)} au repos~;
la fonction \texttt{customer\_priority\_pro\_rata(book, qty, entitlement\_pct)}~; \texttt{price\_improvement\_auction(agency\_qty, stop\_price, side, responses,
book\_customers, one\_other\_pct, many\_others\_pct)} renvoyant le prix final,
les exécutions, la quantité de l'initiateur et le fait que le prix se soit amélioré ou non.
\bfield{Règles} Entiers uniquement~; conservation (les exécutions plus la quantité de
l'initiateur égalent l'ordre)~; les clients ne reçoivent jamais moins que sous toute autre
règle~; la part de l'initiateur au prix final ne dépasse jamais son pourcentage de
l'ordre \emph{original} lorsque d'autres sont présents, et correspond à l'ordre entier lorsque
personne ne répond.
\bfield{Tests d'acceptation} \texttt{code/firm/match/tests/test\_firm\_optmatch.py}~:
les clients d'abord, quel que soit leur temps d'arrivée~; le droit acquis face à la part
au prorata~; les enchères sans répondant, avec un seul et avec plusieurs~; l'amélioration qui
retire l'ordre~; un client vendeur~; la conservation.
\bfield{Pour aller plus loin} Un \omterm{def:m1:options-market-structure:complex}{carnet d'ordres complexes} pour des spreads à deux jambes, avec des prix de
jambe contraints par les carnets simples.
\end{build}

\begin{omsources}
\item Cboe Exchange, rules of the Automated Improvement Mechanism, tel que
déposé auprès de la SEC (SR-CBOE-2019-045 et dépôts ultérieurs, Exhibit 5).
\item US Securities and Exchange Commission, ordonnances sur le Penny Interval
Program, 2020 (par exemple Release 34-89167, Exhibit 5)~; The Options Industry
Council, «~Why certain options trade in penny increments~».
\item OPRA, \emph{Revised OPRA Capacity Projections}, avis aux abonnés aux
données en multidiffusion, 15 septembre 2025.
\end{omsources}

\section{Exercices}

\begin{exercise}[$\star$]\label{exo:m1:options-market-structure:1}
Le niveau de la \cref{fig:m1:options-market-structure:rules} reçoit un ordre
de 12 lots sous \omterm{def:m1:options-market-structure:customer}{priorité client}. Donner les exécutions.
\end{exercise}

\begin{exercise}[$\star$]\label{exo:m1:options-market-structure:2}
Même niveau, 115 lots, \omterm{def:m1:options-market-structure:customer}{priorité client} et un droit de 40~\% pour le
\omterm{def:m1:what-a-trading-firm-does:market-maker}{teneur de marché} désigné. Reproduire à la main les chiffres de la figure.
\end{exercise}

\begin{exercise}[$\star$]\label{exo:m1:options-market-structure:3}
Dans une classe du \omterm{def:m1:options-market-structure:penny}{penny program} qui ne fait pas partie des trois fonds
entièrement en penny, lesquels de ces prix limites sont valides~: 2.43, 2.99, 3.02,
3.05, 12.10, 12.12~? Quel est le spread minimal, en pourcentage du prix, d'une option cotée
autour de 3.10~?
\end{exercise}

\begin{exercise}[$\star\star$]\label{exo:m1:options-market-structure:4}
À partir de l'\cref{dat:m1:options-market-structure:penny}~: donner le taux moyen de messages
par seconde sur une journée de 6.5 heures, le taux de pointe par seconde, et leur
rapport. Donner la bande passante qu'une firme doit provisionner pour recevoir les deux flux avec la
marge de retransmission.
\end{exercise}

\begin{exercise}[$\star\star$]\label{exo:m1:options-market-structure:5}
Un ordre d'agence de 100 lots est mis aux enchères avec un stop de 2.50. Donner
l'exécution de l'initiateur lorsque~: personne ne répond~; une firme répond à 2.50 pour 100~;
trois firmes répondent à 2.50 pour 100 chacune~; une firme répond à 2.49 pour 60 et
deux à 2.50.
\end{exercise}

\begin{exercise}[$\star\star$]\label{exo:m1:options-market-structure:6}
Le call 50 est coté 2.40 à 2.50 et le call 55 0.90 à 1.00. Donner le prix net
de l'achat du spread de calls 50--55 jambe par jambe aux prix affichés, sa
valeur au point médian, et la largeur du marché du package. Pourquoi un \omterm{def:m1:options-market-structure:complex}{carnet d'ordres complexes}
peut-il le coter beaucoup plus serré que 0.20~?
\end{exercise}

\begin{exercise}[$\star\star\star$]\label{exo:m1:options-market-structure:7}
\emph{Programmation.} Avec \texttt{simulate\_auctions(4000, m, 0.5, 24)}, indiquer la part de l'initiateur et la
part de l'ordre ayant reçu un meilleur prix, pour $m = 1$ et $m = 3$ répondants en
moyenne. Qu'est-ce que le client gagne grâce au troisième
répondant~?
\end{exercise}

\begin{exercise}[$\star\star\star$]\label{exo:m1:options-market-structure:8}
\emph{Trouver la faille.} Une société de négociation pour compte propre qui n'est pas enregistrée comme
\omterm{def:m1:what-a-trading-firm-does:market-maker}{teneur de marché}, écrit~: «~Notre cours acheteur de 5~lots a rejoint en premier le meilleur cours
acheteur, devant 450~lots de cotations de \omterm{def:m1:what-a-trading-firm-does:market-maker}{teneurs de marché}. Nous serons exécutés en premier.~» La bourse
applique la \omterm{def:m1:options-market-structure:customer}{priorité client} et l'\omterm{def:m1:matching-algorithms-and-implied-spreads:prorata}{allocation au prorata}. Que se passera-t-il lorsque 100~lots seront
vendus à ce prix~?
\end{exercise}

\section{Problème~: l'enchère}

\begin{problem}[{Problème du week-end --- ce que vaut une part garantie}]\label{pb:m1:options-market-structure:1}
Un grossiste reçoit l'ordre d'un client d'acheter 100~calls. Le marché est à 2.40
contre 2.50~; le grossiste valorise le call à 2.44. Il peut lancer une
\omterm{def:m1:options-market-structure:auction}{enchère d'amélioration de prix} avec un stop à 2.50, selon les règles de
l'\cref{dat:m1:options-market-structure:aim}. D'après l'expérience~: personne ne répond
dans 10~\% des enchères~; une société répond au stop dans 25~\% des cas~; plusieurs
répondent au stop dans 35~\% des cas~; dans 20~\% des cas, les répondants améliorent à 2.49 pour 60~lots, le reste
s'exécutant au stop avec plusieurs participants~; dans 10~\% des cas, les répondants améliorent pour
l'ensemble de l'ordre.

\medskip\noindent\textbf{Partie I --- La part.}
\begin{enumerate}
\item Donner l'exécution du grossiste dans chacun des cinq cas.
\item Donner son exécution espérée.
\item Donner son avantage par lot à 2.50, en dollars par contrat.
\item Donner son profit brut espéré par enchère.
\item Donner l'\omterm{def:m1:retail-flow-and-wholesaling:pfof}{amélioration de prix} espérée pour le client, par contrat.
\end{enumerate}

\medskip\noindent\textbf{Partie II --- Les alternatives.}
\begin{enumerate}[resume]
\item Si le grossiste acheminait simplement l'ordre vers le carnet de la bourse,
combien paierait le client, et combien gagnerait le grossiste~?
\item Le grossiste verse au courtier 30~cents par contrat pour l'ordre.
Donner son profit net espéré par enchère.
\item Il pourrait choisir à la place un stop à 2.48. Qu'est-ce qui change pour le client,
pour l'avantage du grossiste, et pour l'incitation des répondants à améliorer~?
\item Un répondant valorise le call à 2.45. À quels prix répondre vaut-il la peine,
et que lui faut-il pour répondre à temps~?
\end{enumerate}

\medskip\noindent\textbf{Partie III --- Le \omterm{def:m1:what-a-trading-firm-does:market-maker}{teneur de marché} dans le carnet.}
Un \omterm{def:m1:what-a-trading-firm-does:market-maker}{teneur de marché} offrait 50~lots à 2.50 dans le carnet au moment où
l'enchère a débuté.
\begin{enumerate}[resume]
\item Sa cotation au repos participe-t-elle à l'allocation de l'enchère à 2.50~? Que doit-il faire
pour y participer~?
\item Pourquoi pourrait-il préférer ne pas répondre, et garder sa cotation pour les ordres
qui ne sont \emph{pas} mis aux enchères~?
\item Quels ordres ne sont pas mis aux enchères, et qu'est-ce que cela implique pour le flux
qui atteint les cotations affichées~?
\item Comment cela doit-il modifier la largeur de ses cotations affichées~?
\item Relier ceci au \cref{ch:m1:retail-flow-and-wholesaling}.
\end{enumerate}

\medskip\noindent\textbf{Partie IV --- Jugement.}
\begin{enumerate}[resume]
\item Une enchère de 100~millisecondes est-elle concurrentielle~? Donner l'argument dans
chaque sens.
\item Qui paie, au final, le \omterm{def:m1:options-market-structure:marketing}{marketing fee}~?
\item Pourquoi les groupes de bourses exploitent-ils plusieurs bourses d'options
aux règles d'allocation différentes~?
\item Qu'impliquerait un passage à des incréments en centimes dans toutes les classes pour les enchères~?
\item Énoncer le \emph{résultat nommé}~: l'allocation espérée à la société
initiatrice.
\item En une phrase~: que vaut l'ordre d'option d'un client, et pour qui~?
\end{enumerate}
\end{problem}

\section{Questions d'entretien}

\begin{interviewq}[$\star$ \iqroles{trader, chercheur}]\label{iq:m1:options-market-structure:1}
En quoi le marché américain des options diffère-t-il, dans sa structure, du
marché américain des actions~?
\end{interviewq}

\begin{interviewq}[$\star$ \iqroles{trader, développeur}]\label{iq:m1:options-market-structure:2}
Qu'est-ce que la \omterm{def:m1:options-market-structure:customer}{priorité client}, et pourquoi les bourses l'accordent-elles~?
\end{interviewq}

\begin{interviewq}[$\star\star$ \iqroles{trader, chercheur}]\label{iq:m1:options-market-structure:3}
Expliquez-moi une \omterm{def:m1:options-market-structure:auction}{enchère d'amélioration de prix} du point de vue
de l'initiateur et d'un répondant.
\end{interviewq}

\begin{interviewq}[$\star\star$ \iqroles{développeur}]\label{iq:m1:options-market-structure:4}
Vous devez consommer le flux consolidé des options. Quels sont les
problèmes d'ingénierie~?
\end{interviewq}

\begin{interviewq}[$\star\star$ \iqroles{trader, chercheur}]\label{iq:m1:options-market-structure:5}
Pourquoi les \omterm{def:m1:options-market-structure:mm}{teneurs de marché d'options} ont-ils besoin d'une \omterm{def:m1:options-market-structure:protection}{protection
des cotations}, et que coûte-t-elle au reste du marché~?
\end{interviewq}

\begin{interviewq}[$\star\star\star$ \iqroles{chercheur, trader}]\label{iq:m1:options-market-structure:6}
Vous êtes \omterm{def:m1:options-market-structure:mm}{teneur de marché d'options} sur un carnet affiché. En quoi
l'existence des enchères et des grossistes modifie-t-elle le flux que vous voyez, et
que faites-vous à ce sujet~?
\end{interviewq}
